\documentclass[a4paper]{article} %

\date{} %

% Please, do not change any of the above lines

\usepackage{amssymb} % add other LaTeX packages you need here.

\begin{document}
\setcounter{page}{1} % Please, do not delete this line

% Your title
\title{Dual spaces of ortholattices}

% Authors
\author{Andrew P.K. Craig, Miroslav Haviar \\ %
       University of Johannesburg\\Matej Bel University \\ % Affiliation 1
       % New author \\
       % New affiliation \\
       % Add authors and affiliation as needed
       \texttt{miroslav.haviar@umb.sk} % Only one corresponding e-mail, please
       }%

\maketitle

% The abstract

We describe digraphs with topology, which give dual representations of ortholattices. This is done via so-called dual Plo\v{s}\v{c}ica spaces of lattices. We define the dual space of a general ortholattice as the dual Plo\v{s}\v{c}ica space of the lattice reduct of the ortholattice equipped with a map representing the orthocomplement operation. We introduce an abstract ortho-Plo\v{s}\v{c}ica space capturing the properties of the dual space of an ortholattice, and we present dual representation theorems between general ortholattices and the ortho-Plo\v{s}\v{c}ica spaces. We illustrate our dual representations by examples and present an open problem.

%Please, edit it and then upload it as a \LaTeX~file at the link included in the email we sent you.

%The abstract should not exceed two pages, including the bibliography.

\begin{thebibliography}{99}
\bibitem{B84}
L. Beran.
\newblock \emph{Orthomodular lattices. Algebraic Approach,}
\newblock Academia, Prague, 1984.
\bibitem{P3}
A.P.K. Craig, M.J. Gouveia, M. Haviar.
\newblock \emph{TiRS graphs and TiRS frames: a new setting for duals of canonical extensions,}
\newblock Algebra Universalis \textbf{74} 123--138 (2015)
\bibitem{P4}
A.P.K. Craig, M.J. Gouveia, M. Haviar.
\newblock \emph{Canonical extensions of lattices are more than perfect,}
\newblock Algebra Universalis \textbf{83:12} (2022)
\bibitem{P7}
A.P.K. Craig, M. Haviar, K. Marais.
\newblock \emph{Dual digraphs of finite meet-distributive and modular lattices,}
\newblock CUBO \textbf{26} 279--302 (2024)
\bibitem{P1}
A.P.K. Craig, M. Haviar, H.A. Priestley.
\newblock \emph{A fresh perspective on the canonical extensions of bounded lattices,}
\newblock Appl. Categ. Struct. \textbf{21} 725--749 (2013)
\bibitem{P5}
A.P.K. Craig, M. Haviar, J. S\~{a}o Jo\~{a}o.
\newblock \emph{Dual digraphs of finite semidistributive lattices,}
\newblock CUBO \textbf{24} 369--392 (2022)
\bibitem{DOvA}
W. Dzik, E. Orlowska, C. van Alten.
\newblock \emph{Relational Representation Theorems for General Lattices with Negations,}
\newblock Lecture Notes in Computer Science \textbf{4136} 162--176 (2006)
\bibitem{Plos95}
M. Plo\v{s}\v{c}ica.
\newblock \emph{A natural representation of bounded lattices,}
\newblock Tatra Mountains Math. Publ. \textbf{5} 75--88 (1995)
\bibitem{U78}
A. Urquhart.
\newblock \emph{A topological representation theory for lattices,}
\newblock Algebra Universalis \textbf{8} 45--58 (1978)

\end{thebibliography}

\end{document}